% ============================================================
% Thesis — Introduction
% Requires in the preamble: tikz with the positioning, arrows.meta and
% decorations.pathreplacing libraries (loaded in 00_main.tex), plus
% amsmath. All citation keys resolve against mainbib.bib.
% ============================================================

\section{Introduction}
\label{sec:introduction}

\subsection{Motivation}
\label{sec:motivation}

Numerical claims are pervasive in political discourse, news reporting, and social media. They
typically involve quantities, dates, percentages, rankings, or comparisons, and small shifts in scope
or numerical interpretation can change a claim's veracity entirely. This makes numerical claim
verification considerably more demanding than standard textual entailment, because a system must
reason jointly over linguistic context and quantitative relationships rather than over surface
semantics alone~\cite{mishra2022numgluesuitefundamentalchallenging,rahman2025fragilenumbersenseprobing}.
The challenge is amplified in multilingual settings, where the same quantitative assertion may appear
in several languages and where a reliable verifier must preserve numerical meaning across
them~\cite{huang2024survey}.

Large language models (LLMs) have become the standard tool for the reasoning and ranking components of
automated fact-checking, used both as rankers of candidate evidence and as judges of reasoning
quality~\cite{sahitaj2025automatedfactcheckingrealworldclaims,qin2024largelanguagemodelseffective}.
A growing body of work shows, however, that their judgments are not language-invariant: semantically
equivalent inputs may receive different rankings or verdicts depending solely on the language in which
they are expressed~\cite{fu2025reliablemultilingualllmasajudge}. Prompting-based rankers additionally
suffer from systematic presentation biases such as position and verbosity bias, which destabilise
their outputs further~\cite{shi2025judgingjudgessystematicstudy}. These observations sit in tension
with evidence that multilingual LLMs encode partly language-independent internal
representations~\cite{wendler2024llamasworkenglishlatent,STOLLE2026100250}, raising the question of
why such latent consistency does not translate into consistent external behaviour.

This tension motivates the central concern of the thesis. If a model can be encouraged to produce
consistent judgments across languages, does that consistency reflect genuinely more robust reasoning,
or merely more stable surface behaviour with no underlying gain in quality? Put differently: is
language invariance a useful \emph{objective} to optimise for, or only a property worth measuring?

\subsection{Problem Statement and Research Gap}
\label{sec:problem}

The question is sharpened by an observation from the consistency literature that is easy to overlook.
Consistency is formally \emph{decoupled} from
correctness~\cite{elazar2021measuringimprovingconsistencypretrained,Qi_2023}: a model can be
consistently wrong, agreeing with itself across languages on an incorrect answer, or accurate on
average while contradicting itself across translations. Cross-lingual agreement therefore certifies
nothing on its own, and any claim that invariance is beneficial has to be demonstrated against task
quality measured independently---not inferred from the agreement statistics themselves.

Against that backdrop, two strands of recent work bear on the question but have not been brought
together. First, objectives that target cross-lingual consistency directly---consistency
regularisation, cross-lingual self-alignment, and preference optimisation adapted to parallel
prompts~\cite{zheng2021consistencyregularizationcrosslingualfinetuning,wang2025calmunleashingcrosslingualselfaligning,liu2026posttraininglanguagemodelscrosslingual}---show
that invariance can be optimised rather than merely observed. These methods are evaluated on factual
question answering or general benchmarks, however, not on ranking quality in a downstream task, so
whether optimising for invariance \emph{improves the task} remains open. Second, because numerical
magnitude is not explicitly encoded by standard subword tokenisation, LLMs reason unreliably about
quantities~\cite{rahman2025largelanguagemodelsnumberland,wallace2019nlpmodelsknownumbers}, and the
numeral \emph{script} itself is an independent axis of cross-lingual variation: the same quantity
written in Arabic-Indic digits tokenises differently and is handled measurably
worse~\cite{reddy2026effectscriptsformatsllm}. Value-aware numerical representations have been
proposed to address this~\cite{dutulescu2026valueawarenumericalrepresentationstransformer}, but
whether they improve cross-lingual robustness rather than only monolingual numeracy is unestablished.

There is therefore no systematic account of how a language-invariance objective, and improved
numerical representations, interact in multilingual numerical fact-checking. This thesis addresses
that gap directly, treating language invariance not only as a property to be measured but as an
explicit target that a training objective can be designed to promote---and then asking, with accuracy
and consistency measured separately, what optimising for it actually buys.

\subsection{Grounding: CLEF CheckThat! 2026 Task 2}
\label{sec:grounding}

We study these questions in the concrete setting of CLEF~CheckThat!~2026 Task~2 on multilingual
numerical claim verification~\cite{clef-checkthat:2026:task2}. The task adopts a test-time scaling
formulation for fact verification~\cite{chungkham2025thinkrightmoretesttime}: for each claim and its
retrieved evidence, multiple reasoning traces are generated by an LLM under varying temperatures to
induce diverse reasoning paths, redundant traces are removed, and participants must train a verifier
that ranks the remaining traces by their utility for reaching the correct verdict and then derives a
claim-level verdict from the top-ranked traces. The English claims come from the QuanTemp
benchmark~\cite{v2024quantemprealworldopendomainbenchmark}, together with Spanish and Arabic claims
from the previous CLEF edition. Ranking quality is measured with Recall@5 and MRR@5 over traces whose
verdict matches the gold label, claim-verdict prediction with macro-averaged F1, and the per-language
score is the mean of the two.

This setting suits the research question for three reasons. It is genuinely multilingual and evaluated
per language, so cross-lingual behaviour is directly observable. Its intermediate structure---an
explicit ranking over reasoning traces---exposes not just \emph{what} a system concludes but
\emph{which evidence} it relied on, which turns out to be where the interesting cross-lingual effects
live. And because the claims are numerical, correctness is grounded in quantities that are invariant
across language and phrasing~\cite{lewkowycz2022solvingquantitativereasoningproblems}, making
inconsistency attributable to the model rather than to genuine semantic drift between translations.

The thesis builds on our own participation in the shared task~\cite{khawaja2026gipplab}. That system
fine-tunes Llama-3.1-8B-Instruct~\cite{llama31instruct} with Quantized Low-Rank
Adaptation~\cite{dettmers2023qloraefficientfinetuningquantized,hu2021loralowrankadaptationlarge} as a
multilingual supervised trace scorer, scoring each claim--trace pair independently and ranking traces
by predicted relevance. This formulation proved more controllable than generative ranking and ranked
first on English and third on Spanish on the official leaderboard, providing a competitive baseline
and an established pipeline on which the present work builds.

\begin{figure}[t]
\centering
% The diagram is naturally wider than \textwidth; scale it to the text block
% so its edges line up with the surrounding paragraphs.
\resizebox{\linewidth}{!}{%
\begin{tikzpicture}[
    >=Stealth,
    font=\small,
    claimbox/.style={draw, rounded corners, fill=blue!8, minimum width=2.2cm,
        minimum height=1.0cm, align=center},
    langbox/.style={draw, rounded corners, fill=gray!10, minimum width=1.9cm,
        minimum height=0.8cm, align=center},
    modelbox/.style={draw, rounded corners, fill=orange!15, minimum width=2.2cm,
        minimum height=1.7cm, align=center},
    outbox/.style={draw, rounded corners, fill=green!10, minimum width=2.5cm,
        minimum height=0.8cm, align=center},
    note/.style={align=center, font=\footnotesize\itshape}
]
\node[claimbox] (claim) {Numerical\\ claim $c$\\ (+ evidence)};
\node[langbox, right=0.9cm of claim, yshift=1.4cm]  (en) {English $c_{\mathrm{en}}$};
\node[langbox, right=0.9cm of claim]                (es) {Spanish $c_{\mathrm{es}}$};
\node[langbox, right=0.9cm of claim, yshift=-1.4cm] (ar) {Arabic $c_{\mathrm{ar}}$};
\node[modelbox, right=0.9cm of es] (model) {LLM\\ trace scorer\\ $f_\theta$};
\node[outbox, right=0.9cm of model, yshift=1.4cm]  (oen) {$r_{\mathrm{en}}$, $\hat{y}_{\mathrm{en}}$};
\node[outbox, right=0.9cm of model]                (oes) {$r_{\mathrm{es}}$, $\hat{y}_{\mathrm{es}}$};
\node[outbox, right=0.9cm of model, yshift=-1.4cm] (oar) {$r_{\mathrm{ar}}$, $\hat{y}_{\mathrm{ar}}$};
\draw[->] (claim) -- (en);
\draw[->] (claim) -- (es);
\draw[->] (claim) -- (ar);
\draw[->] (en) -- (model);
\draw[->] (es) -- (model);
\draw[->] (ar) -- (model);
\draw[->] (model) -- (oen);
\draw[->] (model) -- (oes);
\draw[->] (model) -- (oar);
\draw[decorate, decoration={brace, amplitude=6pt}]
    (oen.north east) -- (oar.south east);
\node[note, right=0.5cm of oes, text width=2.7cm] {%
    \textbf{Invariance:} an ideal model yields consistent rankings $r$ and
    verdicts $\hat{y}$ across all three language versions of $c$};
\end{tikzpicture}%
}
\caption{The language-invariance problem in multilingual numerical fact-checking. The same numerical
claim $c$ and its evidence are presented to the trace scorer $f_\theta$ in three languages. A
language-variant model may produce divergent rankings $r$ and verdicts $\hat{y}$ for semantically
equivalent inputs, whereas a language-invariant model should produce consistent outputs. This thesis
asks whether explicitly optimising for such invariance improves ranking and reasoning quality.}
\label{fig:intro-language-invariance}
\end{figure}

Figure~\ref{fig:intro-language-invariance} illustrates the core problem. The same claim, together with
its evidence and candidate reasoning traces, is presented to the trace scorer in English, Spanish, and
Arabic. A language-variant model can return divergent rankings and verdicts for these semantically
equivalent inputs, whereas a language-invariant model should return consistent ones. Quantifying and
reducing this divergence---and determining whether reducing it also improves ranking
quality---is the empirical heart of this thesis.

\subsection{Research Questions}
\label{sec:rqs}

Guided by the gap described above, this thesis investigates the relationship between language
invariance and reasoning quality in multilingual numerical fact-checking through three research
questions:

\begin{itemize}
    \item[\textbf{RQ1}] Does optimising for language-invariant reasoning improve ranking performance
    in multilingual numerical fact-checking?
    \item[\textbf{RQ2}] What is the relationship between language invariance and reasoning quality
    across different languages?
    \item[\textbf{RQ3}] Can language-invariant training objectives lead to more robust multilingual
    reasoning models?
\end{itemize}

\noindent
Cutting across these questions is a methodological sub-thread on representation: to what extent
value-aware numerical representations, which expose quantitative magnitude more explicitly than
ordinary subword tokenisation, contribute to cross-lingual robustness in numerically grounded
reasoning. We treat this as a lens on RQ2 and RQ3 rather than as a separate question, since numerical
grounding is precisely where language invariance is most testable.

\subsection{Contributions}
\label{sec:contributions}

In addressing these questions, this thesis makes the following contributions:

\begin{enumerate}
    \item It formulates language invariance as an explicit, differentiable \emph{training objective}
    rather than only an evaluation criterion, defined at the level of per-trace relevance scores: the
    objective penalises a model when the normalised score profiles it assigns to the same reasoning
    traces diverge across languages. It operationalises this on a competitive trace-scoring pipeline
    derived from our CLEF~CheckThat!~2026 Task~2 system.
    \item It provides a systematic comparison of ranking formulations through the lens of
    cross-lingual consistency, spanning zero-shot listwise and pairwise prompted ranking, Direct
    Preference Optimization~\cite{rafailov2024directpreferenceoptimizationlanguage}, and supervised
    trace scoring, and shows why a pointwise scorer is the formulation that makes invariance
    optimisable at all.
    \item It contributes an evaluation protocol that reports task quality and cross-lingual agreement
    \emph{independently}, and at two levels---the final verdict and the underlying trace
    ranking---and demonstrates that these two levels give substantially different answers.
    \item It reports an empirical study over three open-weight backbones and $28$ non-degenerate
    configurations, evaluated on a held-out trilingual test set of $4{,}233$ parallel claims with
    three seeds per configuration, including the interaction between the invariance objective and
    value-aware numerical representations.
\end{enumerate}

\paragraph{Summary of findings.}
The headline result is a separation rather than an improvement. Optimising for language invariance
leaves the official ranking metric essentially unchanged---Recall@5 is statistically indistinguishable
from standard multilingual fine-tuning, and both sit near the metric's attainable ceiling---while
improving the cross-lingual \emph{stability} of the ranking roughly fourfold, with the
invariance-trained and non-invariance populations not overlapping at all on that measure. A modest
gain in verdict quality accompanies it, concentrated in Arabic and specifically in the minority
\textsf{Conflicting} class. Two secondary findings support the case: verdict-level agreement is high
for every system we evaluated and therefore fails to discriminate between them, so a protocol that
measures consistency only over final answers would have found nothing; and a trivially invariant
random baseline scores perfectly on every agreement measure while being useless, which is the sharpest
possible demonstration that consistency must be reported alongside, never instead of, task quality.

\subsection{Thesis Outline}
\label{sec:outline}

The remainder of this thesis is organised as follows.
\textbf{Section~\ref{sec:related-work}} reviews related work on LLM-based ranking and its presentation
biases, automated and multilingual fact-checking, language invariance and cross-lingual consistency,
cross-lingual representations, numerical reasoning across languages, and the preference-optimisation
and parameter-efficient fine-tuning methods on which the training objectives build.
\textbf{Section~\ref{sec:background}} introduces the concepts and terminology the rest of the thesis
assumes, for readers unfamiliar with language models, ranking formulations, fact-checking pipelines,
or the agreement statistics used to quantify consistency.
\textbf{Section~\ref{sec:methodology}} presents the methodology: the task formulation and data,
the trace-scoring architecture and its training objective, the score-level language-invariance
objective, the value-aware numerical components, the baselines, and the evaluation protocol.
\textbf{Section~\ref{sec:results}} reports and discusses the results, presenting task performance and
cross-lingual consistency separately before combining them to answer the research questions, and
states the threats to validity.
\textbf{Section~\ref{sec:conclusion}} synthesises the findings, revisits the contributions, and
outlines directions for future work.
Full per-model results are provided in Appendix~\ref{app:results-full}.
